% Draft: §6, What each analysis states, for the restructured paper (rewrite
% method agreed 17 Oct 2026; see the evidence repo's DECISIONS.md). Drafted from
% the old draft's §9 (its comparison table and §9.5), STATE-OF-THE-ARGUMENT
% items 6--11 and 14, SECTION-NOTES ("§4 passes to §6"; "§2 passes to §6 and
% §8"), and the §3 and §4 drafts, whose pages it reuses. New page: Payne et al.
% p. 574, property (iv), checked against the PDF on 17 Oct. Revised after the
% muse review (reviews/muse-sec6-20261017-221323/TRIAGE.md).
% Version A of ordinary Head throughout (Brett, 17 Oct). The old §9's rows on
% phrase category and modifier domains, Det:{NP,PP}, the subsection on degree
% modifiers and the shared-projection condition of failure are dropped.
% Pending: the online supplement has to be rebuilt on version A before the
% pointer below is true. The row for once, twice and thrice was added 18 Oct
% with §5 (drafts/once.tex), and extended after the §5 muse review, with a
% term-of-office clause below the line (reviews/muse-sec5-20261018-101536/TRIAGE.md).
% 18 Oct: the pre-head adjective row, the trade paragraph and the last
% paragraph now point to §4's accounting of pronouns and names (CGEL pp. 423,
% 430). Not yet muse-reviewed.
%
% Label: sec:costs (nothing refers to it yet). Relies on sec:analyses,
% sec:det-head, sec:mod-head, sec:fused-relatives.

\section{What each analysis states}\label{sec:costs}

Sections~\ref{sec:det-head} and~\ref{sec:mod-head} compared the two analyses construction by construction. Table~\ref{tab:costs} puts the results together. Both columns use \textit{CGEL}'s categories and the same lexical facts, so what the table compares is how each analysis has to state them. The full fragment, with the permissions and Head relations of each analysis, is in the online supplement.

\begin{table}[H]
\centering\footnotesize
\caption{What fusion and ordinary Head state for determinatives standing alone and for adjectives heading a nominal. Above the line, the two analyses state the same content, each in its own terms; below it, they differ.}\label{tab:costs}
\setlength{\tabcolsep}{4pt}\renewcommand{\arraystretch}{1.2}
\begin{tabular}{>{\raggedright\arraybackslash}p{3.9cm}>{\raggedright\arraybackslash}p{5.5cm}>{\raggedright\arraybackslash}p{5.5cm}}
\toprule
 & Fusion (\textit{CGEL}) & Ordinary Head \\
\midrule
Determinatives that stand alone (\mention{some}, normally not \mention{every}) & A list of the words that allow Det--Head & The same list, as permissions to head a nominal with no determiner \\
The forms \mention{none} and \mention{mine} & Dependent form as determiner, independent form as head (p.~410) & The same \\
Determinatives that head a nominal after a determiner (\mention{the few}) & A default from the modifier use, with exceptions both ways: head but not modifier (\mention{some}, \mention{both}), modifier but not head (\mention{every}, \mention{various}) & The same default and exceptions, as permissions \\
A post-head modifier in that use, cardinals apart (\mention{the many who did}) & Normally needed (p.~416) & The same condition \\
Determiners these heads take & Word by word: a plural noun's range for \mention{few} (\mention{those chosen few}); for adjectives, \mention{the} or a genitive with \mention{rich} (p.~417), none with \mention{blue} & The same \\
\mention{A few} and its expansions (\mention{a lucky few}) & A slot for a modifier after \mention{a} in the DP headed by \mention{few} (Section~\ref{sec:mod-head}), or a list of sequences with \mention{a} that take a plural head (p.~353) & The same \\
Adjectives that head a nominal & \textit{CGEL}'s types, with exclusions such as \ungram{\mention{the main}} (pp.~415--417) & The same list, as permissions \\
Genitive on a determinative standing alone & The gap predicted for all, with a list of exceptions (\mention{either's}, \mention{neither's}, \mention{each's}; compounds such as \mention{everyone's}, unless their nominal base carries it) & The genitive expected, with a list of gaps (\mention{all}, \mention{many}, \mention{few}, \mention{some}; \mention{once}, \mention{twice}, \mention{thrice}) \\
Change of category on the head path & Where a DP or AdjP heads a nominal, or a PP (\mention{up to twice}, Payne et al., p.~589 n.~31) & The same \\
\mention{Once}, \mention{twice}, and \mention{thrice} & Numerical uses fused like \mention{one}: Det--Head on their own, Mod--Head after a determiner; singulative and past uses Det--Head only (Payne et al., pp.~588, 592); term of office largely confined to names of offices (p.~593); sequential use not fused (a preposition, p.~594) & The same uses, as permissions: numerical uses head a nominal with or without a determiner, singulative and past uses only without one; the same range for term of office; sequential use as for fusion \\
\midrule
A second function for the head of the nominal & Determiner for a determinative with no determiner before it; modifier for one after a determiner, and for adjectives and nouns such as \mention{nylon}; term-of-office \mention{once} a nominal of its own as attribute (Payne et al., p.~593) & None outside fused relatives, so \mention{few} heads its nominal in the same way in \mention{few would dare}, \mention{few of them}, \mention{the few}, and \mention{the lucky few}; term-of-office \mention{once} a modifier, as \mention{two} is in \mention{the two books} \\
A pre-head adjective with no determiner (\ungram{\mention{Lucky few escaped}}) & Excluded by adjacency, as for the compounds (p.~423), but \mention{Poor everybody!} needs an exception (Section~\ref{sec:mod-head}) & Excluded by a condition: only after a determiner. \textit{CGEL} states a restriction of this kind for pronouns (p.~430), and \mention{everybody} and \mention{everyone} join their exception \\
\bottomrule
\end{tabular}
\end{table}

Above the line, the content is shared. The lists, defaults, and conditions follow the word, as Sections~\ref{sec:det-head} and~\ref{sec:mod-head} found, and each analysis phrases them in its own terms: fusion as licences for a fused function, ordinary Head as permissions to head a nominal.

Below the line is the trade. Fusion gives the head of the nominal a second function and derives from it the absence of a pre-head adjective before a bare determinative, with an exception for \mention{Poor everybody!}. Ordinary Head has no second function and states the condition instead, a statement of the kind \textit{CGEL} already makes for pronouns (p.~430; Section~\ref{sec:mod-head}). It gains generality: \mention{few} heading the same kind of NP gets one structure, where fusion gives it one pair of functions with no determiner and another after one. That's the kind of gain \textit{CGEL} claims for fusion against traditional grammar (p.~419; Section~\ref{sec:mod-head}).

The two departures \textcite[570--571]{Payne2007} distinguish fall differently too. Fusion puts both on one constituent wherever a DP, AdjP, or PP is a fused head, since that phrase has two functions and is also where the category changes. Ordinary Head never puts both on one constituent. Its constituents with two functions are in fused relatives, where the fused phrase and its parent must project the same lexical category (p.~574), so no category changes there. In a bare fused relative such as \mention{I'll take \bk{whichever he gave me}}, the two meet in one NP at different nodes: the DP heads the nominal, and the NP is head and prenucleus at once (Section~\ref{sec:fused-relatives}, Figure~\ref{fig:fusedrelthis}).

On the criterion of Section~\ref{sec:analyses}, a property of form that varied with the determiner or modifier function and not with the word would count for fusion. Form there means the shape of words, order, and the dependents and determiners a word takes. For the determiner function, the test compares one word with and without it. \textit{CGEL} gives determinatives no use as a pure head, so the comparison is between a determinative standing alone and the same word after a determiner, and it changes two things at once: the word loses its determiner function, and a determiner precedes it. A fact that separates the two uses is predicted by fusion, and also by any condition that mentions a determiner. The modifier function has no such contrast, and the test is the one Section~\ref{sec:mod-head} applied: a form found only with the function, or a fact of order that only a modifier relation predicts. Shape, dependents, the determiners taken, and the genitive show nothing of either kind (Sections~\ref{sec:det-head} and~\ref{sec:mod-head}).

The pre-head adjective does, and it counts for fusion. \mention{Few} takes none standing alone in edited prose, and takes one after \mention{the}, and fusion derives the difference from adjacency, where ordinary Head states a condition. \mention{A lucky few} isn't an exception on either analysis of \mention{a} (Section~\ref{sec:mod-head}). \mention{Poor everybody!} is, as \mention{Lucky you!} is for \textit{CGEL}'s statement about pronouns. So the evidence for the determiner function is one fact with one exception, and against it are the generality ordinary Head gains for \mention{few} and the constituent with two parents that fusion gives every NP headed by a bare determinative. I take the balance to favour ordinary Head, though not to force it. What would reverse it is a second fact of the adjective's kind: a form a determinative has only when it stands alone, which its meaning and countability don't explain.

The saving is in where fusion applies, not in whether the grammar has it. Ordinary Head keeps constituents with two functions for fused relatives, so it removes them from noun phrases without a relative clause, not from the grammar. What it gains there is one head relation for each word across its uses, and constituents with two parents confined to the one construction where form shows the second function. Section~\ref{sec:fused-relatives} finds such facts in fused relatives.
